\scanpage{170}

\uline{Corollaire 1} \deleted{Pour que $f$ soit définie en une \unsure{la} sous-variété irréductible $Y$, il faut et il suffit que $Y$ ne soit \unsure{contenu|contenue} dans aucun \unsure{des} \unsure{$V$} tels que} \inserted{L'ens. des pts de $X$ où $f$ n'est pas définie est la réunion des supports des $V$ tels que $v_V(f)<0$.} C'est \unsure{encore vrai} si $f$ est une application rationnelle de $X$ dans une variété \deleted{\uline{affine}} ($X$ normal) c'est par suite une sous-variété pure de dim $n-1$.

\deleted{\unsure{Th.}} N.B. Les résultats précédents valent sans changement essentiel si \deleted{\unsure{X n'est pas}} \inserted{l'ens} des pts singuliers de $X$ est de dim $\leq n-2$.

\uline{Th.} Soit $X$ une variété \deleted{non singulière}\inserted{normale de dim $n$}, $Y$ \deleted{\unsure{q}} une variété complète, $f:X\longrightarrow Y$ une application rationnelle. Alors l'ens. des pts où $f$ n'est pas définie est de dim $\leq n-2$.

\deleted{On peut supposer $X$} Deux démonstrations

a) On peut supposer, en regardant le \uline{graphe} de $f$, que $f$ est inverse d'un \uline{morphisme} birationnel $g:Y\longrightarrow X$, et $g$ est \deleted{propre} \uline{propre}.

\deleted{L'ens} Soit $X'$ l'ens des pts $x$ \deleted{\unsure{de}} tels que $g^{-1}(x)$ soit de dim $>0$, alors $X'$ est de dim $\leq n-2$, car sinon $g^{-1}(X')$ serait de dim $\geq(n-1)+1=n$, absurde. \deleted{Donc $g^{-1}(X)$} D'autre part, il résulte du Main th. qu'au dessus de $\complement X'$, $g$ est un isom. birégulier, donc $g^{-1}$ est définie sur $\complement X'$. O.K.

N.B. Bien entendu, le th. se généralise aussi au cas où \deleted{\unsure{$X$}} l'ens des singularités de $X$ est de dim $\leq n-2$

% archive sheet number: 169
